% Główny plik dokumentu – kompiluj właśnie ten plik

\documentclass[a4paper,12pt,oneside,onecolumn]{book}
% Ustawienia formatu dokumentu:
% - a4paper: format A4
% - 12pt: wielkość czcionki
% - oneside: jednostronny druk
% - onecolumn: jedna kolumna tekstu

\linespread{1.5} % Zwiększenie interlinii (1.5x) dla lepszej czytelności
% Kodowanie i czcionka
\usepackage[utf8]{inputenc}         % Kodowanie UTF-8
\usepackage[T1]{fontenc}            % Lepsze dzielenie wyrazów i znaki diakrytyczne
\usepackage{helvet}                 % Czcionka Helvetica
\renewcommand{\familydefault}{\sfdefault} % Domyślna czcionka bezszeryfowa

% Matematyka i symbole
\usepackage{amsmath,amsfonts,amssymb,amsthm} % Rozszerzenia matematyczne


% Marginesy
%\usepackage[margin=2.5cm]{geometry} % Ustawienie marginesów
\usepackage[top=2.5cm, bottom=2.5cm, left=3.5cm, right=1.5cm]{geometry} % lewy i prawwy inaczej dla marginesow 


% Ulepszenia typografii
\usepackage{microtype}
\usepackage{xcolor}
\usepackage{array}
\usepackage{enumitem} % opisy do wzorow
% ODSTEPY ENUMERATE 1. 2. 3. 
\setlist[enumerate]{itemsep=0pt, parsep=0pt, topsep=0pt}

% Odnośniki i linki
\usepackage{url}
\usepackage[hidelinks]{hyperref}   % Linki bez kolorów/ramki
\urlstyle{same} % linki taka sama czcionka



% --- FORMATOWANIE NAGŁÓWKÓW (Musi być dodane) ---
\usepackage{titlesec}

% Formatowanie ROZDZIAŁÓW (Nagłówek 1)
% Wymogi: 16 pkt, pogrubiona, odstępy 18pkt przed i po
\titleformat{\chapter}[block]
  {\bfseries\fontsize{16}{16}\selectfont} % Styl: Pogrubienie, rozmiar 16pt
  {\chaptername\ \thechapter.} % Numeracja: "Rozdział X."
  {0.5em} % Odstęp między "Rozdział X." a tytułem
  {}

% Odstępy dla rozdziału: {lewy}{góra}{dół}
\titlespacing*{\chapter}{0pt}{-20pt}{18pt}

% Formatowanie ROZDZIAŁÓW nienumerowanych (Spisy, Literatura, Wstęp)
\titleformat{name=\chapter,numberless}[block]
  {\bfseries\fontsize{16}{16}\selectfont}
  {}
  {0pt}
  {}
\titlespacing*{name=\chapter,numberless}{0pt}{-20pt}{18pt}

% Formatowanie PODROZDZIAŁÓW (Nagłówek 2 - np. 1.1)
% Wymogi: 14 pkt, pogrubiona, odstępy 12pkt przed i po
\titleformat{\section}[block]
  {\bfseries\fontsize{14}{14}\selectfont}
  {\thesection.}
  {0.5em}
  {}
\titlespacing*{\section}{0pt}{12pt}{12pt}

% Formatowanie PODPUNKTÓW (Nagłówek 3 - np. 1.1.1)
% Wymogi: 12 pkt, pogrubiona, odstępy 6pkt przed i po
\titleformat{\subsection}[block]
  {\bfseries\fontsize{12}{12}\selectfont}
  {\thesubsection.}
  {0.5em}
  {}
\titlespacing*{\subsection}{0pt}{6pt}{6pt}



% Wcięcia i bibliografia ************************************************************************************
\usepackage{indentfirst}           % Wcięcie w pierwszym akapicie każdego rozdziału
\setlength{\parindent}{1cm}

% --- BIBLIOGRAFIA I CYTOWANIA (Poprawiona) ---
\usepackage[polish]{babel}
\usepackage{csquotes}
\MakeOuterQuote{"}

% Konfiguracja BibLaTeX - styl authoryear (APA-like)
\usepackage[
    style=authoryear,    % Styl: Autor (Rok)
    backend=biber,       % Silnik przetwarzania
    maxcitenames=2,      % W tekście: max 2 nazwiska, potem "i in." (zgodnie z APA)
    maxbibnames=99,      % W bibliografii: wymień wszystkich autorów
    uniquelist=false,    % Nie dodawaj inicjałów, jeśli nazwiska się powtarzają
    isbn=false,          % WYŁĄCZ ISBN (zgodnie z szablonem)
    doi=false,           % WYŁĄCZ DOI (zgodnie z szablonem)
    url=false,           % WYŁĄCZ URL dla książek/artykułów
    eprint=false
]{biblatex}

% Włączamy URL tylko dla typu @online (netografia)
\AtEveryBibitem{%
  \ifentrytype{online}{}{\clearfield{url}}%
  \ifentrytype{online}{}{\clearfield{urldate}}%
}

% --- FORMATOWANIE WYGLĄDU bibliografi ---

% 1. Wszystkie tytuły kursywą (książki, artykuły, strony www)
\DeclareFieldFormat{title}{\mkbibemph{#1}}          
\DeclareFieldFormat[article]{title}{\mkbibemph{#1}} 
\DeclareFieldFormat[book]{title}{\mkbibemph{#1}}    
\DeclareFieldFormat[online]{title}{\mkbibemph{#1}}    

% 2. USUWANIE "W:" przed nazwą czasopisma
% Standardowo biblatex daje "In:" lub "W:". Ta komenda to usuwa.
\renewbibmacro{in:}{}

% 3. Kolejność: Wydawnictwo, Miasto (Dla książek)
% Standard to Miasto: Wydawnictwo. UE Katowice chce odwrotnie.
\renewbibmacro*{publisher+location+date}{%
  \printlist{publisher}%
  \setunit*{\addcomma\space}%
  \printlist{location}%
  \setunit*{\addcomma\space}%
  \usebibmacro{date}%
  \newunit}

% 4. Separatory
\renewcommand*{\nameyeardelim}{\addcomma\space} % Przecinek po nazwisku w cytowaniu
\renewcommand*{\postnotedelim}{\addcomma\space} % Przecinek przed numerem strony

% 5. Data dostępu dla netografii: [dostęp DD.MM.RRRR]
\DefineBibliographyStrings{polish}{
  urlseen = {dostęp},
}
\DeclareFieldFormat{urldate}{\mkbibbrackets{\bibstring{urlseen}\space#1}}
\DeclareFieldFormat{url}{\url{#1}}

% 6. LISTA NUMEROWANA (1. 2. 3.) do bibliografi 
\newcounter{mybibcounter}
\defbibenvironment{bibliography}
  {\list
     {\stepcounter{mybibcounter}\themybibcounter.} 
     {\setlength{\leftmargin}{2.5em}%              
      \setlength{\labelwidth}{2em}%                
      \setlength{\labelsep}{0.5em}%                
      \setlength{\itemindent}{0pt}%
      \setlength{\itemsep}{\bibitemsep}%
      \setlength{\parsep}{\bibparsep}}}
  {\endlist}
  {\item}





% Kod źródłowy (np. programy)
\usepackage{listings}

% Nagłówki i stopki
\usepackage{fancyhdr}
\usepackage{etoolbox}

% Grafika
\usepackage{graphicx}
\graphicspath{ {images/} }         % Folder z grafikami

% Opisy pod rysunkami
\usepackage{caption}
\usepackage{chngcntr}

% Globalne numerowanie rysunków i tabel
\counterwithout{figure}{chapter}
\counterwithout{table}{chapter}

%definicja podpisu obrazka
\DeclareCaptionFont{elevenpt}{\fontsize{11}{13.6}\selectfont} % Definicja 11pkt
\captionsetup{
  font=elevenpt, % Użycie 11pkt
  labelfont=bf,
  labelsep=period,
  justification=raggedright, 
  singlelinecheck=false
}


\usepackage{float}

\newcommand{\MyFigure}[5][0.8\textwidth]{%
  \begin{figure}[H] % H dla tego miejsca ht dla ustawienia roznie 
    \centering
    \includegraphics[width=#1]{#2}
    \caption{#3}
    \label{#4}
    \vspace{-3pt}
    {\raggedright\fontsize{10pt}{12pt}\selectfont\textit{Źródło: #5}\par}
  \end{figure}
}

% Niestandardowe środowisko do tabel z podpisem i źródłem ---- TABELE
\captionsetup[table]{
    format=plain,
    justification=raggedright,
   labelfont=bf,
   textfont=normalfont,
   labelsep=period,
   name=Tabela,
   singlelinecheck=false,
   skip=2pt
}

 \newcommand{\MyTableSource}[1]{%
   \vspace{-8pt}
   \begin{flushleft}
     {\linespread{0.9}\fontsize{10}{12pt}\selectfont\parbox{0.95\textwidth}{\raggedright\textit{Źródło: #1}}}
   \end{flushleft}
   \vspace{-6pt}
 }

% Nadpisanie stylu rozdziału – aby używał fancyhdr
\patchcmd{\chapter}{\thispagestyle{plain}}{\thispagestyle{fancy}}{}{}

% Stylowanie kodu źródłowego
\definecolor{codegreen}{rgb}{0,0.6,0}
\definecolor{codegray}{rgb}{0.5,0.5,0.5}
\definecolor{codepurple}{rgb}{0.58,0,0.82}
\definecolor{backcolour}{rgb}{0.98,0.98,0.95} % Nieco jaśniejsze tło

\lstdefinestyle{mystyle}{
    backgroundcolor=\color{backcolour},
    commentstyle=\color{codegreen},
    keywordstyle=\color{blue}\bfseries, % Słowa kluczowe na niebiesko i pogrubione
    numberstyle=\tiny\color{codegray},
    stringstyle=\color{codepurple},
    basicstyle=\ttfamily\footnotesize,
    breakatwhitespace=false,
    breaklines=true,
    captionpos=b,
    keepspaces=true,
    numbers=left,
    numbersep=10pt,                  % Większy odstęp numerów od kodu
    showspaces=false,
    showstringspaces=false,
    showtabs=false,
    tabsize=4,                       % Standard dla Pythona to 4 spacje
    frame=single,                    % Ramka wokół kodu
    rulecolor=\color{black!10},      % Kolor ramki
    extendedchars=true,              % Obsługa znaków specjalnych
    literate={ą}{{\k{a}}}1 {ć}{{\'c}}1 {ę}{{\k{e}}}1 {ł}{{\l{}}}1 {ń}{{\'n}}1 {ó}{{\'o}}1 {ś}{{\'s}}1 {ź}{{\'z}}1 {ż}{{\.z}}1 {Ą}{{\k{A}}}1 {Ć}{{\'C}}1 {Ę}{{\k{E}}}1 {Ł}{{\L{}}}1 {Ń}{{\'N}}1 {Ó}{{\'O}}1 {Ś}{{\'S}}1 {Ź}{{\'Z}}1 {Ż}{{\.Z}}1 % Polskie znaki w komentarzach
}

% Definicja konkretnie pod Pythona
\lstdefinestyle{pythonstyle}{
    style=mystyle,
    language=Python,
    morekeywords={self, cls, @property, @staticmethod, yield}
}

\lstset{style=pythonstyle} % Ustawienie jako domyślny

% Spolszczenia nazw elementów
\renewcommand{\lstlistingname}{Listing}
\renewcommand{\lstlistlistingname}{Spis listingów}
\renewcommand{\listfigurename}{Spis rysunków}
\renewcommand{\listtablename}{Spis tabel}
\renewcommand{\figurename}{Rysunek}
\renewcommand{\tablename}{Tabela}
\renewcommand{\chaptername}{Rozdział}
\renewcommand{\contentsname}{Spis Treści}

% Środowisko do warunków (np. w matematyce)
\newenvironment{conditions}
  {\par\vspace{\abovedisplayskip}\noindent\begin{tabular}{>{$}l<{$} @{${}={}$} l}}
  {\end{tabular}\par\vspace{\belowdisplayskip}}



% --- NOWE ŚRODOWISKO "GDZIE" (opis do wzorow) ---
\newenvironment{gdzie}
{
    \vspace{-1.0em}    % Podciągnij tekst do góry (bliżej wzoru)
    \noindent \textbf{gdzie:}  % Automatyczny napis "gdzie:"
    \begin{itemize}[         % Lista z ciasnymi ustawieniami
        noitemsep,           % Brak odstępów między punktami
        topsep=2pt,  % Mały odstęp od "gdzie:" do pierwszego punktu
        parsep=0pt,          % Brak odstępów wewnątrz punktów
        partopsep=0pt,       % Brak dodatkowych odstępów
        leftmargin=1.5em,    % Wcięcie listy
        labelsep=0.5em     % Odstęp kropki/myślnika od tekstu
    ]
}
{
    \end{itemize}
    \vspace{0.5em}                  % Odstęp po zakończeniu opisu
} % Dołączenie pliku z definicją stylu dokumentu

% Dołączenie plików z bibliografią (literatura i netografia)
\addbibresource{refs.bib}
\addbibresource{url_refs.bib}


\usepackage[titles]{tocloft}


\begin{document}

% *************** Front matter ***************
% -*- TeX-master: "./main.tex" -*-
% *************** Strona tytułowa ***************

\pagestyle{empty}
\noindent
\begin{center}
    UNIWERSYTET EKONOMICZNY W KATOWICACH
\end{center}

\begin{center}
    INFORMATYKA
\end{center}

\vfill
\begin{center}
    \textbf{Łukasz Osadnik}
    \par \normalsize \textbf{157249}
\end{center}
\vfill

\begin{center}
    \Large
	\textbf{Budzik z wykorzystaniem mikrokontrolera}
\end{center}

\begin{center}
	\textbf{Alarm clock utilizing microcontroller}
\end{center}

\vfill
\vfill
\vfill

\begin{flushright}
Praca licencjacka \\
pod kierunkiem Doktora Tomasza Staśa
\end{flushright}

\vfill
\begin{center}
    KATOWICE 2026
\end{center}

\newpage 

\pagestyle{fancy}
\fancyhf{}
\fancyhf{} % sets both header and footer to nothing
\renewcommand{\headrulewidth}{0pt}
\fancyfoot[R]{\thepage}

\begin{center}
    \textbf{OŚWIADCZENIE PROMOTORA}
\end{center}

\par \noindent Oświadczam, że niniejsza praca została przygotowana pod moim kierunkiem
i stwierdzam, że spełnia wymogi stawiane pracom dyplomowym. \\ 
\par \noindent Pracę akceptuję. \\ \\ 

\begin{minipage}{2in}
\begin{center}
\ldots\ldots\ldots\ldots\ldots\ldots\ldots\ldots\ldots\ldots{} \\
(data)
\end{center}
\end{minipage}
\hfill
\begin{minipage}{2in}
\begin{center}
\ldots\ldots\ldots\ldots\ldots\ldots\ldots\ldots\ldots\ldots{} \\
(podpis promotora)
\end{center}
\end{minipage}

\newpage

% *************** Oświadczenie początkowe ***************
\begin{minipage}[t]{0.5\textwidth}
\begin{flushleft}
\ldots\ldots\ldots\ldots\ldots\ldots\ldots\ldots\ldots\ldots{} \\
Imię i nazwisko

\ldots\ldots\ldots\ldots\ldots\ldots\ldots\ldots\ldots\ldots{} \\
Kierunek

\ldots\ldots\ldots\ldots\ldots\ldots\ldots\ldots\ldots\ldots{} \\
Nr albumu
\end{flushleft}
\end{minipage}
\hfill
\begin{minipage}[t]{0.5\textwidth}
 \begin{flushright}
Katowice, dnia \ldots\ldots\ldots\ldots\ldots\ldots{} 
 \end{flushright}
\end{minipage}

\begin{center}
    \textbf{OŚWIADCZENIE}
\end{center}

\par \noindent Mając świadomość odpowiedzialności prawnej oświadczam, że złożona praca licencjacka pt. Budzik z wykorzystaniem mikrokontrolera\\
została napisana przeze mnie samodzielnie.
Równocześnie oświadczam, że praca ta nie narusza praw autorskich w rozumieniu ustawy z dnia 4 lutego 1994 roku o prawie autorskim i prawach pokrewnych oraz dóbr osobistych chronionych prawem.
Ponadto praca nie zawiera informacji i danych uzyskanych w sposób niedozwolony i nie była wcześniej przedmiotem innych procedur związanych z uzyskaniem dyplomów lub tytułów zawodowych uczelni wyższej.
Wyrażam zgodę na nieodpłatne udostępnienie mojej pracy w celu oceny jej oryginalności przez Jednolity System Antyplagiatowy prowadzony przez ministra właściwego ds. szkolnictwa wyższego oraz przechowywania jej w Ogólnopolskim Repozytorium Prac Dyplomowych, oraz wewnętrznej bazie prac dyplomowych Uniwersytetu Ekonomicznego w Katowicach. Oświadczam, że poinformowano mnie o zasadach dotyczących oceny oryginalności pracy dyplomowej przez Jednolity System Antyplagiatowy.
Oświadczam także, że ostateczna wersja pracy przesłana przeze mnie drogą elektroniczną jest zgodna z plikiem poddanym ocenie w Jednolitym Systemie Antyplagiatowym.
Jednocześnie oświadczam, że jest mi znany przepis art. 233 § 1 Kodeksu karnego określający odpowiedzialność za składanie fałszywych zeznań.

\begin{flushright}
\ldots\ldots\ldots\ldots\ldots\ldots\ldots\ldots\ldots\ldots\ldots\ldots\ldots{} \\
(podpis składającego oświadczenie)
\end{flushright}


% *************** Spis treści ***************
\tableofcontents

% *************** Koniec front matter ***************

 % Tutaj znajduje się strona tytułowa, spis treści, oświadczenia itd.

% *************** Treść rozdziałów *******-> 80-130 tys znakow
\chapter*{Wstęp}
\addcontentsline{toc}{chapter}{Wstęp}
Wprowadzenie do tematu, omówienie czym jest budzik oraz mikro kontroler.
\chapter{Projekty innych ludzi}
Przedstawienie projektów innych ludzi.
Na chwilę obecną na google scolar udało mi się znaleść.
\section{Zegarki}
Choć mój projekt nie ma na celu prezentacji czasu na ekranie lcd lub w podobny sposób.
Ważną częścią tego projektu jest pilnowanie która jest godzina.
\newline\newline
\parencite{lukman2016development} oraz \parencite{khan2012designing}
\section{Budziki}
\parencite{screws470digital}
\section{Inne}
Inne projekty które mają wspólne cech z tym co chciałem osiągnąć.

\chapter{Co chcę osiągnąć}
Wstęp do tego co chcemy osiągnąć
\section{Cele osobiste}
Celem jest zlikwidowanie potrzeby korzystanie z telefony jako budzik'a
\section{Cele ogólne}
Chciałbym stworzyć alarm z wykorzystaniem mikrokontroler'a ukrytego w pluszaku który można sterować na różne sposoby z telefonu lub komputer'a.
Dodatkowo chciał bym udostępnić możliwość tworzenia własnych dźwięków budzika użytkowniką.
\chapter*{Teoria}
\addcontentsline{toc}{chapter}{Teoria}
krótki wstęp do tego co chcemy zrobić
\chapter{Technologie}
krótki wstęp do wykorzystywanych technologii do osiągniecia naszego celu
\section{Raspberry pi pico 2}
mikro kontroler który będę programować
\subsection{C++}
język który będę wykorzystywać
\subsection{FreeRTOS}
bibliteka/system do zadań asynchronicznych
\subsection{Pi pico maker mini}
nakładka (HAT) raspberry pi pico który będę wykożystywać
\newline\newline
posiada złącze na baterie oraz bzyczek
\section{Opakowanie}
opis metod wykorzystywanych do tworzenie pluszaka w którym zamkniemy nasz mikrontroler
\section{Klient do alarmu}
(z angielskiego client) interfejsy po których komunikacja z mikrokontrolerem będzie możliwa
\subsection{Strona internetowa}
strona wykożystujące interfejs REST API mikrokontrolera do komunikacji z nim
\newline\newline
dodatkowo kożystająca z technologii PWA/serviceworkerów pozwalającej wczytać pliki graficzne, html, css oraz javascript raz a następnie przechowywanie ich w przeglądarce.
\subsection{CLI}
tak zwany commandline interface pozwalał by na komunikacje z urządzeniem z terminal'a również za pomocą REST API
\newline\newline
napisany albo w C++ albo Deno/Typescript (Łatwiejsza opcja).
\subsubsection{Projekt interfejsu CLI}
projekt faktycznego interfejsu wierszu poleceń zgodnie z filozowią unix (programów robięcych jedną rzecz dobrze które można ze sobą połączyć)
\subsection{Aplikacja natywna}
napisana w bibliotece GTK4/Adwait'a (możliwa zmiana) pozwalała by na kontrole alarmu bez użycia przeglądarki lub terminala zarówno na telefonie jak i komputerze.
\newline\newline
dodatkowo dzięki wykorzystaniu natywnej biblioteki zamiast pakowania Electro'u aplikacja będzie znacznie miej zasobożerna niż przeglądarka internetowa.
\chapter*{Praktyka}
\addcontentsline{toc}{chapter}{Praktyka}
wstęp do faktycznej budowy gotowego rozwiązania
\chapter{Oprogramowanie mikrokontrolera}
w tej sekcji zajmować się będę tworzeniem oprogramowania do mikrokontrolera.
\section{Komunikacja bezprzewodowa}
programowanie chip'u wifi obecnego na mikrokontrolerze
w tym API REST
\section{NTP}
programowanie systemu pilnowania czasu
\section{Bzyczek}
programowanie sygnału dźwiękowego
\chapter{Opakowanie dla mikrokontrolera}
wstęp do budowy opakowanie
\section{Tworzenie pluszaka}
proces tworzenia pluszaka w którym zamkniemy nasz mikrokontroler
\section{Wypełnianie wnętrza pluszaka}
składanie peryfieriów i kabli do kupy.
\chapter{Tworzenie strony internetowe}
wstęp do tworzenia strony internetowej
\section{Budowa interfejsu webowego}
w tej sekcji opisuwany będzie proces tworzenia strony internetowej do interakcji z urządzeniem.
\section{PWA}
wykorzystanie technologii pwa do zapisywania strony w przeglądarce internetowej.
\chapter{Budowa aplikacji CLI}
wstęp do tworzenia aplikacji cli
\section{Implementacji wcześniej zaprojektowanego interfejsu}
w tej sekcji opisuwany będzie proces tworzenia aplikacji cli.
\chapter{Programowanie aplikacji natywnej}
wstęp do tworzenia aplikacji natywnej
\section{Tworzenie interfejsu graficznego}
w tej sekcji opisuwany będzie proces tworzenia aplikacji natywnej od strony wyglądu.
\section{Programowanie interfejsu graficznego}
w tej sekcji będziemy programować funkcjonalną część aplikacji
\chapter*{Podsumowanie}
\addcontentsline{toc}{chapter}{Podsumowanie}
Podsumowanie wyników mojej pracy
To jest mój projekt i się udał.
\chapter{Testy i wnioski}
Raport z faktycznego użytku przygotowanego rozwiązania
\chapter{Co dalej?}
co można by było ulepszyć/dodać do mojego projektu

% % -*- TeX-master: "../main.tex" -*-
% Wstęp
\chapter*{Wstęp}
\addcontentsline{toc}{chapter}{Wstęp}

Lorem ipsum dolor sit amet, consectetur adipiscing elit. Integer vel diam nec sapien tincidunt sodales. Curabitur non nisi vel sapien dignissim efficitur. Nulla facilisi. Sed id arcu mi. Vestibulum viverra ex vel nisl suscipit fermentum. Lorem ipsum dolor sit amet, consectetur adipiscing elit. Integer vel diam nec sapien tincidunt sodales. Curabitur non nisi vel sapien dignissim efficitur. Nulla facilisi. Sed id arcu mi. Vestibulum viverra ex vel nisl suscipit fermentum.


% \chapter{Przykładowy rozdział z ilustracjami, tabelą i tekstem}
% Rozpoczęcie nowego rozdziału – jego tytuł pojawi się w spisie treści i nagłówkach

\section{Wprowadzenie}
% Sekcja – tytuł pojawi się w spisie treści

% Przykładowy tekst akapitowy. 
% Możesz tu dodać wprowadzenie teoretyczne, kontekst badawczy, przegląd literatury itd.
% Zastosowanie \cite{} służy do cytowania źródeł z pliku refs.bib (literatura).
Lorem ipsum dolor sit amet, consectetur adipiscing elit. Integer vel diam nec sapien tincidunt sodales. Curabitur non nisi vel sapien dignissim efficitur. Nulla facilisi. Sed id arcu mi. Vestibulum viverra ex vel nisl suscipit fermentum. Lorem ipsum dolor sit amet, consectetur adipiscing elit. Integer vel diam nec sapien tincidunt sodales. Curabitur non nisi vel sapien dignissim efficitur. Nulla facilisi. Sed id arcu mi. Vestibulum viverra ex vel nisl suscipit fermentum. Lorem ipsum dolor sit amet, consectetur adipiscing elit. Integer vel diam nec sapien tincidunt sodales. Curabitur non nisi vel sapien dignissim efficitur. Nulla facilisi. Sed id arcu mi. Vestibulum viverra ex vel nisl suscipit fermentum. Lorem ipsum dolor sit amet, consectetur adipiscing elit. Integer vel diam nec sapien tincidunt sodales. Curabitur non nisi vel sapien dignissim efficitur. Nulla facilisi. Sed id arcu mi. Vestibulum viverra ex vel nisl suscipit fermentum \parencite{PortalAI_CzymJestSI}.

\section{Wzory matematyczne}
% Sekcja poświęcona wzorom matematycznym – używa środowiska equation

% Poniżej przykład prostego wzoru matematycznego
\begin{align}
    f(x) = ax^2 + bx + c
    \label{eq:funkcja_kwadratowa} % Etykieta umożliwia odwołanie się do wzoru w tekście
\end{align}

\begin{gdzie}
    \item $a$ -- to cos 
    \item $b$ -- to cos 
    \item $x$ -- to cos 
\end{gdzie}


% Odwołanie do wzoru przez \ref{}
% Przykład: "Zob. równanie (1)"
% Równanie \ref{eq:funkcja_kwadratowa} ...
Równanie \ref{eq:funkcja_kwadratowa} może mieć dwa, jeden lub zero pierwiastków rzeczywistych w zależności od wartości wyróżnika trójmianu kwadratowego.

\section{Analiza i interpretacja}
% Sekcja główna z dwiema podsekcjami

\subsection{Podrozdział pierwszy}
% Podsekcja z dodatkowym opisem

% Przykładowy długi akapit, w którym można opisać analizę, interpretacje, wyniki itd.
Lorem ipsum dolor sit amet, consectetur adipiscing elit. Proin euismod, sapien nec feugiat fermentum, eros purus pulvinar augue, nec pretium nulla nunc sed neque. Sed ut turpis in felis tincidunt efficitur Lorem ipsum dolor sit amet, consectetur adipiscing elit. Proin euismod, sapien nec feugiat fermentum, eros purus pulvinar augue, nec pretium nulla nunc sed neque. Sed ut turpis in felis tincidunt efficitur Lorem ipsum dolor sit amet, consectetur adipiscing elit. Proin euismod, sapien nec feugiat fermentum, eros purus pulvinar augue, nec pretium nulla nunc sed neque. Sed ut turpis in felis tincidunt efficitur Lorem ipsum dolor sit amet, consectetur adipiscing elit. Proin euismod, sapien nec feugiat fermentum, eros purus pulvinar augue, nec pretium nulla nunc sed neque. Sed ut turpis in felis tincidunt efficitur \parencite{article}.
% \parencite{} – cytowanie w nawiasie (APA)

\subsection{Podrozdział drugi z tabelą}
% Przykład dodania tabeli do dokumentu

% Środowisko table z komendą tabular – tabela z trzema kolumnami
\begin{table}[ht]
\centering
\caption{Przykładowa tabela danych}
\label{tab:przykladowa_tabela}
\begin{tabular}{|>{\raggedright\arraybackslash}p{4cm}|c|c|}
\hline
\textbf{Kategoria} & \textbf{Wartość A} & \textbf{Wartość B} \\
\hline
Lorem ipsum        & 12                & 34                 \\
Dolor sit amet     & 45                & 67                 \\
Consectetur        & 89                & 10                 \\
\hline
\end{tabular}

\vspace{4pt} % Mały odstęp między tabelą a źródłem
{\raggedright\fontsize{10pt}{12pt}\selectfont\textit{Źródło: Opracowanie własne}\par}
\end{table}



% Odwołanie do tabeli
Dane przedstawione w tabeli \ref{tab:przykladowa_tabela} mają charakter poglądowy.

\section{Wizualizacja}
% Sekcja z rysunkiem/grafiką

% Dodanie rysunku przy użyciu zdefiniowanego wcześniej polecenia \MyFigure
% Parametry:
% - ścieżka do pliku
% - podpis
% - etykieta (do odwołania)
% - źródło (np. publikacja, własna ilustracja itd.)
\MyFigure{images/example-image.jpg}{Przykładowy rysunek}{rys:przykladowy_rysunek}{\parencite{article}}



% Odwołanie do rysunku
(rysunek \ref{rys:przykladowy_rysunek}) jest ilustracją poglądową, pokazującą użycie komendy \texttt{\textbackslash MyFigure}.


\section{Numeracja}

\begin{enumerate}
    \item item1

    \item item2
\end{enumerate}

\section{Implementacja kodu }
\begin{lstlisting}[style=pythonstyle, caption={Przykład zastosowania }, label={code:przyklad}]
import os

# pogladowy kod
X = 2+1+1+5
\end{lstlisting}

\clearpage
% Wymuszenie rozpoczęcia nowej strony (opcjonalne)

% % -*- TeX-master: "../main.tex" -*-
\chapter{Przykładowy rozdział z ilustracjami, tabelą i tekstem}

\section{Podsumowanie}

Lorem ipsum dolor sit amet, consectetur adipiscing elit. Integer vel diam nec sapien tincidunt sodales. Curabitur non nisi vel sapien dignissim efficitur. Nulla facilisi. Sed id arcu mi. Vestibulum viverra ex vel nisl suscipit fermentum. Lorem ipsum dolor sit amet, consectetur adipiscing elit. Integer vel diam nec sapien tincidunt sodales. Curabitur non nisi vel sapien dignissim efficitur. Nulla facilisi. Sed id arcu mi. Vestibulum viverra ex vel nisl suscipit fermentum. Lorem ipsum dolor sit amet, consectetur adipiscing elit. Integer vel diam nec sapien tincidunt sodales. Curabitur non nisi vel sapien dignissim efficitur. Nulla facilisi. Sed id arcu mi. Vestibulum viverra ex vel nisl suscipit fermentum. Lorem ipsum dolor sit amet, consectetur adipiscing elit. Integer vel diam nec sapien tincidunt sodales. Curabitur non nisi vel sapien dignissim efficitur. Nulla facilisi. Sed id arcu mi. Vestibulum viverra ex vel nisl suscipit fermentum.


% \input{chapters/ending}


% *************** Bibliografia ***************
% wymagania - Praca musi zawierać co najmniej 25 pozycji literaturowych, w tym 4 pozycje obcojęzyczne. Przypisy muszą być zgodne ze stylem APA. jako {nazwisko, rok, strona}
\newpage
\chapter*{Literatura} % Rozdział bez numeru
\addcontentsline{toc}{chapter}{Literatura} % Dodanie do spisu treści

\section*{Bibliografia}
\setcounter{mybibcounter}{0} 
\printbibliography[heading=none, notkeyword=netografia] % Wyświetla źródła nieoznaczone jako netografia

\section*{Netografia}
\setcounter{mybibcounter}{0} 
\printbibliography[heading=none, keyword=netografia] % Wyświetla źródła oznaczone jako netografia


% *************** Spisy pomocnicze ***************
\clearpage
\phantomsection
\addcontentsline{toc}{chapter}{Spis rysunków}
\listoffigures

\clearpage
\phantomsection
\addcontentsline{toc}{chapter}{Spis tabel}
\listoftables

% Jeśli chcesz dodać spis kodów (próbek kodu) 
\clearpage
\phantomsection
\addcontentsline{toc}{chapter}{Spis listingów}
\lstlistoflistings

\end{document}


